\section{A direct unbounded proof for the literal 238-state table}
\label{sec:runtime}

The positive block argument above explains the state saving and provides
one proof route. A second route executes the exact new table symbolically
from every guarded source instruction frame, with no bound on represented
register values. This section states the frame precisely and explains
why the finite proof construction gives an unbounded theorem.

\subsection{Frames, the logical seed and spare storage}

Use inherited labels in this section, so dispatcher 96 is canonical state
84. A frame is in that dispatcher with the head at the first pointer
cell and a one-marker immediately to its left. The finite pointer word
is followed by zero padding until a header one at distance $64+g$ from
the head, where $g\in\mathbb N$ is an arbitrary additional gap. After the
header come nine primary register blocks and a nonempty finite list of
spare blocks. A block of value $v$ is
\[
                       B(v)=0\,1^{v+1}.
\]
The tape outside this complete finite frame is zero. This convention
keeps the separators, the extra unary one in every block, the final
spare tail and the zero exterior in the specification.

The logical-to-physical register permutation is
\[
   (\operatorname{phy}(0),\ldots,\operatorname{phy}(8))
                      =(6,2,0,4,3,7,8,5,1).
\]
The initial logical register function and its physical representation are
\begin{align*}
 R_*&=(3,0,0,0,3,0,0,0,0),& R_*(i)&=0\quad(i\geq9),\\
 R_{\rm physical}&=(0,0,0,3,0,0,3,0,0).
\end{align*}
The 167 initial spare blocks each have value one. They are auxiliary
binary storage, not the logical register function above index eight.
The source only addresses registers zero through eight, so a logical
tail of zeros is preserved throughout source execution.

Two different tape implementations execute the new table from the
all-zero tape: a sparse set of one-cells in inherited labels and a dense
array in canonical labels. Both first reach the dispatcher after exactly
409,835 transitions with head $-5$, marker $-6$, header $176$, and gap
$g=117$. They agree on every tape cell and decode the displayed primary
values and 167 spare ones. There are 351 one-cells, ending at position
701; every other cell is zero. The complete position list is retained in
\file{evidence/full_replay/bootstrap.json}. The independently checked
configuration-by-configuration positive simulation supplies the same
complete frame from the boundary initializer.

\subsection{The instruction theorem}

\begin{proposition}[Whole-tape instruction segments]
\label{prop:runtime}
For each of the 6,455 occurrences of $P_1$, every natural primary register
assignment satisfying the selected decrement guard, every natural gap,
and every nonempty finite spare list, the literal new table executes a
finite positive segment implementing exactly that source instruction.
For a nonhalting instruction it returns to a complete valid dispatcher
frame with the prescribed next pointer and updated registers. The head
position and entire tape agree with the specified output. For the unique
halt instruction it reaches the designated halt state. No nonhalting
case or intermediate portion is accepted as a halt.
\end{proposition}

There are 1,931 decrement occurrences. Checking both outcomes for each,
and one case for every other occurrence, yields $6455+1931=8386$ cases.
A positive decrement starts its tested value at $1+x$ with $x\in\mathbb N$;
its zero case fixes only that value to zero. All other primary values
remain independent natural parameters. The cases describe all values,
not a finite sample of small registers.

An increment adds one to the selected register and appends a spare block
of value one. A positive decrement subtracts one and increases the
additional gap by one. Breaks, continues and zero decrements perform
the exact pointer operations stated in the source section. Pointer
overflow can shift the head and marker one cell left while adjusting the
gap. The comparison checks the actual motion separately from the tape
representation, so agreement cannot be obtained merely by translating
the output frame. Every nonhalting endpoint retains a nonempty spare
list.

\subsection{A finite language for infinite families of tapes}

The inherited proof engine stores the tape as a zipper: a finite symbolic
word to each side of the head, an orientation and a control state. Word
atoms are constant bits, runs whose lengths are affine forms, and an
arbitrary finite list of unary blocks. Its thirteen affine slots are a
constant, the extra gap, nine primary values, the first spare value,
and the total width of the remaining spare list. The last width is the
actual width of that list, not a separate unconstrained parameter that
may disagree with it. Reversal flags retain list orientation as the head
crosses between zipper sides.

The representation also uses a second block family
\[
                        C(v)=1\,0\,1^v.
\]
All affine coefficients are integers. Where a literal bit must exist,
the proof demands a positive constant term and nonnegative coefficients.
Where an entire variable run is traversed, a zero length is permitted and
handled by the induction base case. This difference prevents reading a
symbol from a possibly empty symbolic run.

\subsection{The sound inference rules}

The proof construction has three execution rules and exact word
normalization. Their soundness follows from ordinary induction.

\paragraph{Literal transition.}
If a run is known nonempty, apply the actual table row to its first bit,
record the write and one-cell motion, and add one to the duration. This
changes exactly one concrete tape cell under every parameter assignment.

\paragraph{Unary self-loop.}
When the actual table row stays in the same state and moves in the scan
direction, it can be applied to an arbitrary nonnegative run length $n$
at once. Induction on $n$ proves the accumulated word, duration $n$ and
motion $\pm n$. For $n=0$, this is the identity computation; it is not an
extra transition. The rule checks the literal row instead of assuming
that a state name denotes a scan loop.

\paragraph{Finite-list traversal.}
Before any list acceleration, the engine derives four one-block
identities from the supplied table, with an arbitrary natural value $v$
and opaque exterior tape. In inherited state labels these are:
\begin{center}
\begin{tabular}{ccccc}
\toprule
State & Direction & Input block & Output block & Steps / absolute motion \\
\midrule
73 & $R$ & $B(v)$ & $B(v)$ & $v+2$ \\
95 & $R$ & $B(v)$ & $C(v)$ & $v+2$ \\
48 & $L$ & $B(v)$ & $B(v)$ & $v+2$ \\
112 & $L$ & $C(v)$ & $B(v)$ & $v+2$ \\
\bottomrule
\end{tabular}
\end{center}
The exit state equals the entry state, permitting induction over any
finite list. The empty list gives the identity; the cons step combines
the one-block result with the induction hypothesis. Total duration and
motion are the sum of the block widths. The one-block proofs use literal
steps and the preceding unary induction, not the list rule they justify.

\paragraph{Exact word identities.}
Adjacent equal-symbol runs combine by addition; zero runs disappear;
reversal retains the denoted word. For a finite value list $A$, one also
uses
\[
                         C(A)\,1=1\,B(A)
\]
and its reversal, proved by base/cons induction from the one-block
identity. This equation explains the apparently shifted boundary one
when a converted spare list is normalized. The proof manipulates whole
words; it does not discard a damaged boundary cell.

These rules preserve a concrete interpretation for every choice of the
parameters and arbitrary spare tail. The engine records both a symbolic
duration and an independent head-motion form. A successful nonhalt case
must have exactly the target control state, the entire normalized left
and right tape words, and the target motion. Its duration must have
nonnegative coefficients and a strictly positive constant. A halt case
can finish only at the actual halt state. None of the accelerated loops
or block identities enters halt in its interior.

\subsection{The completed universal certificates}

For the direct 238-state proof, every one of the 8,386 cases succeeds.
The replay constructs 2,255,513 symbolic rule applications and covers all
476 literal transition rows. The least constant term of a completed
instruction duration is 13. The fresh boundary-239 guard proof succeeds
on the same 8,386 cases with 2,247,513 rules, covering all 478 old rows,
and explicitly verifies the extra neighbor guard as described earlier.

\begin{center}
\begin{tabular}{lrr}
\toprule
Metric & Direct 238 proof & Boundary-239 guard proof \\
\midrule
Guarded cases & 8,386 & 8,386 \\
Symbolic rule applications & 2,255,513 & 2,247,513 \\
Literal rows covered & 476 & 478 \\
Minimum duration constant & 13 & 13 \\
Explicit old $169/1$ runtime visits & Not required & 32 \\
Minimum guarded one-run constant & Not required & 23 \\
\bottomrule
\end{tabular}
\end{center}

The engine's 4,000-rule per-case cap limits the finite attempt to build
a derivation. It puts no bound on a natural parameter, the number or
values of spare blocks, or the runtime represented by an accepted
derivation. Every list and unary acceleration has a proved induction
rule, so the resulting certificates establish Proposition~\ref{prop:runtime}
for unbounded data.

The engine file is unchanged from the attributed coordinated-reuse
package and has SHA-256
\begin{center}\small\ttfamily
9ff1b7c487036744206976aa45c16c3b8\\
d6eeb8d4e7cad81f41dae8ebab66c5e.
\end{center}
Its two wrappers use the new exact tables and pin the boundary source.
They share this proof engine; the word ``separate'' does not mean two
independent implementations of universal symbolic execution.

\subsection{Concatenation and exact source behavior}

At each nonhalting endpoint, the source-control certificate verifies
that the successor pointer and the reached leaf codeword have the same
infinite zero-padded word. Thus the next symbolic instruction theorem
applies to the very same tape, including overflow and shorter-pointer
cases. Translation of a frame along $\mathbb Z$ preserves all literal
steps, so the theorem applies at the actual head position reached by the
previous segment. Induction supplies an exact binary segment for every
finite source prefix starting at $(\epsilon,R_*)$.

Every segment has positive duration. Hence infinitely many source steps
produce an unbounded number of binary steps; there is no infinite
zero-time stuttering. Conversely, every represented nonhalting source
instruction finishes in finite time, so binary nontermination cannot be
caused by a hidden infinite loop within an instruction. The unique halt
occurrence reaches the unique designated binary halt. Together with the
verified blank prefix, this proves
\[
 \Halt(\Mnew,\mathrm{blank})
            \longleftrightarrow \RegHalts(P_1,(\epsilon,R_*)).
\]
